\chapter{Introdução}
\label{cap:introducao}

% \textbf{Introdução -- Convite -- Exibição}

% São as marcas iniciais e ainda inarticuladas do texto, ancoradas em marcas bem articuladas para o leitor ao qual se dirige;
%     o convite se concentra nas contingências do outro, é um kit de odores, um artifício que sobrepõe uma série bem articulada ao seu público com outra inarticulada, carregando uma nova sensibilidade a ser elaborada como um novo órgão dos corpos leitores:
%         qual o nome deste órgão?

% Este órgão carrega uma outra razão, reticulativa, metamórfica;
%     um cérebro de camaleão, o grito do morcego:
%         é um órgão que situa, desenvolve para reenvolver --- a máscara de SAL.

% Aqui o texto-rito será apresentado nos termos da pesquisa científica usual, ainda que subvertendo-os;
%     uma pergunta de pesquisa orienta o texto -- como falar cientificamente? -- e constitui sua intenção;
%         o método, por sua vez -- a heterotopografia -- foi (será) inventado por mim como uma espécie de transposição da etnografia;
%             os resultados e sua discussão são o próprio texto-rito;
%                 o convite tateia uma descrição pseudo-objetivista dessas coisas, mas todas se entrelaçam numa só, que só se conhece lendo tudo;
    
% O fio de Ariadne será a controvérsia em torno de Nego Bispo, a ser heterocartografada no texto-rito de modo a informar os monoteístas da ciência moderna.

% ---

% As contingências dos Outro são as coordenadas diretrizes da escrita;
%     a máscara do convite é a língua do convidado:
%         a língua do anfitrião, neste caso, é o objeto do convite, o qual deve funcionar como artifício primeiro para a articulação-com aquela.

% Serão três introduções -- mancha, máscara e marca -- que apresentam nos termos da língua acadêmica atual o texto-rito que iniciam e ao qual convidam o leitor;
%     uma pergunta de pesquisa, um método, resultados e discussões:
%         utilizarei a heterotopografia para sugerir como falar cientificamente através do próprio texto-rito que também se autodiscute.

% Controvérsias com Nego Bispo darão o fio de Ariadne.

% ---

% {\large \textbf{Mancha}}

% \textit{Objetivo:} especular e descrever um mundo que não seja inadequado \textit{ser assim}.\\[-3pt]
% \hspace*{2cm}\textit{ser assim:} habitar cada momento sem apagar as contingências.

% \textit{Pergunta de pesquisa:} como falar cientificamente?\\[-3pt]
% \hspace*{2cm}\textit{Hipótese/Proposição:} a prática científica orientada pela comunicação científica.\\[-3pt]
% \hspace*{3cm}\textit{Contra-objetivismo} = Nego Bispo + ficar com o problema (da objetividade)\\[-3pt]
% \hspace*{4cm}\textit{Monoteísmo x Politeísmo}

% \textit{Método:} heterotopografia.\\[-3pt]
% \hspace*{2cm} . signo x totem
% \hspace*{2cm} . nomeação como instauração de sinergia heterotópica-semântica
% \hspace*{2cm} . especulação orientada à proposição\\

% \textbf{\textit{A ciência alçada ao patamar da poesia/magia.}}\\

% {\large \textbf{Máscara}}

% baleias, mundo comum (Latour), ir até o fim para saber do que se trata, o problema: o subjetivismo, Givânia do quilombo de Conceição das Crioulas -> pacto feminino com a santa, localização postergada e diferencial, objetivo x PFSD, o pacto contra-objetivista como baleia para uma nova bússola para a ciência -> politeísmo (Nego Bispo), texto espiralar\\

% {\large \textbf{Marca}}

% \vspace{1cm}\noindent
% ---\\
% ---\\
% ---\\

\indent

\encruza{A MANCHA}

Primeira introdução, mais tradicional, na qual antecipo e localizo o meu trabalho de pesquisa;
    um convite orientado às contingências do Outro, especialmente sua língua, de modo a lhe ser envolvente:
        exposição de pergunta de pesquisa, "objeto", arcabouço teórico, "objetivos", métodos, hipóteses e (in)conclusões.

{\small
\begin{description}
    \item[Pergunta de Pesquisa] Como falar cientificamente?
    \item["Objeto"] A controvérsia em torno de Nego Bispo.
    \item["Objetivos"] Compor uma gramática científica capaz de lidar com controvérsias, contradições, paradoxos e incompletudes.
    \item[Arcabouço teórico] Latour, Haraway, Mauss, Foucault, Leda, Krenak, Rufino, Manuela Carneiro (Cultura com aspas).
    \item[Métodos] Heterotopografia, artificiaria.
    \item[Hipóteses] Proposições, Contra-objetivismo, Totem.
    \item[(In)conclusões] Cecinepá, um kit-de-odores para o mundo-comum.
\end{description}}

\encruza{A MÁSCARA}

Articulação entre a realização do texto-rito e as coordenadas contingentes da ciência moderna contemporânea;
    justificativas:
        exposição e desenvolvimento dos emaranhados entre os componentes apresentados na Mancha.

O meu contexto pessoal impulsionando a busca por uma linguagem para se falar melhor sobre autismo;
    o contexto global de crise climática:
        a intenção de elaborar uma nova forma de se falar cientificamente como mediação para a composição contínua de um mundo-comum.

A polêmica em torno do pensamento de Nego Bispo;

Meu objetivo, ou melhor dizendo, minha \textit{intenção} 

\encruza{A MÁSCARA}